\section{Related Work}
\label{sec:related}

\para{Activation verbalizers.}
Training-free methods patch a hidden state into a prompt of the same model and let it describe what it sees~\citep{ghandeharioun2024patchscopes,chen2024selfie}. LatentQA~\citep{pan2024latentqa} trains a decoder, a LoRA copy of the target, to answer open-ended questions about activations. Its questions mostly concern qualitative properties of the coming reply, such as a hidden persona or goal. Activation Oracles~\citep{karvonen2025activation} scale this recipe to about one million examples and release checkpoints for \qwen. Follow-up work finds that narrowly trained AOs match but do not exceed linear probes, and that a text-only control does about as well~\citep{bauer2026building}. AOs also default to ``No'' on yes/no questions and are best scored by answer likelihood~\citep{bauer2026building,torrielli2026confidence}. Other studies report concept-specific blind spots~\citep{bersia2026blind}, hallucination on harmful-intent inputs~\citep{zhao2026faithful}, and strong effects of which token positions are read~\citep{torrielli2026selecting}. None of these works asks about the \emph{timing} of a future event. We evaluate an off-the-shelf AO on exactly this question and train AO-style readers on future-event labels.

\para{Future information in hidden states.}
Future Lens decodes tokens a few steps ahead from a single hidden state~\citep{pal2023futurelens}. \citet{wu2024planahead} argue that most future-useful information consists of ``breadcrumbs'' that are useful now. \citet{dong2025emergent} probe prompt activations for global attributes of the response, such as length and the number of reasoning steps, and find that verbalized self-prediction is near chance. \citet{merzouk2026remaining} show that linear probes encode the number of remaining output tokens. These works motivate our probes and our timing head. Unlike them, we target a rare safety event hundreds of tokens away, under censoring.

\para{Streaming safety monitors.}
Streaming content monitors stop harmful outputs early by classifying partial responses~\citep{li2025judgment,sharma2025constitutional}. These monitors detect harm that is already present in the prefix. StreamGuard~\citep{kavumba2026predict} instead forecasts risk, using Monte Carlo rollouts judged by a guard model as targets, but it reads text only. \citet{qi2025shallow} show that safety alignment is concentrated in the first few tokens, so a compliant opening can lead to harm later. \citet{karny2026surfacing} surface activation-based risk across multi-turn conversations. We focus on lead distance within a single response and compare activation readers against text readers on the same prefixes.

\para{Safety representations and self-report.}
Harmfulness and refusal are linearly readable from residual streams~\citep{zou2023representation,arditi2024refusal}, and such representations can be used for control~\citep{zou2024circuit}. Models can also predict some of their own behavior~\citep{binder2024looking}. We use linear and MLP probes as the main representation baselines and ask \qwen itself whether its response will turn harmful. Our prompts come from HarmBench~\citep{mazeika2024harmbench}, AdvBench~\citep{zou2023universal}, StrongREJECT~\citep{souly2024strongreject}, BeaverTails~\citep{ji2023beavertails}, and Alpaca~\citep{taori2023alpaca}.
